\begin{abstract}
The goal-frontier maximization (GFM) sequence establishes a computable
    alignment objective with structural safety properties and endogenous
    anti-monopolar pressure. The preceding papers treat agent evaluation
    as a statistical signal-processing problem: trust-weighted votes,
    correlational contraction detection, max-aggregated risk reports.
    This paper introduces a structural causal model for capability
    dynamics that upgrades scorpion detection and risk evaluation
    from statistical to causal identification, and sharpens the
    remaining agent-evaluation mechanisms, under explicitly stated
    structural assumptions.
    Do-calculus contraction attribution sharpens scorpion
    detection under additive-SCM, causal-sufficiency, and
    exogenous-orthogonality conditions.
    Pathway-conditional residuals formalize risk-trust with $L^2$
    convergence guarantees. A log-linear opinion pool replaces
    max-aggregation without discarding independence information. A
    probabilistic dependency-risk model sharpens the minimax substrate
    bound into a per-step expected cost that compounds with growth,
    gated by an adversarial balance condition. The common thread:
    a GFM actor needs a causal model of other agents, and the SCM
    provides the structural backbone that each mechanism draws on.
\end{abstract}
